\medskip

On considère la fonction $f$ définie sur $\R$ par 

\[f(x) = \dfrac34 x^2 - 2x + 3.\]

\begin{enumerate}
\item Dresser le tableau de variations de $f$ sur $\R$.
\item En déduire, que pour tout $x$ appartenant à l'intervalle $\left[\dfrac43~;~2\right]$,\: $f(x)$ appartient à l'intervalle $\left[\dfrac43~;~2\right]$.
\item Démontrer que pour tout $x$ réel, $x \leqslant f(x)$.

Pour cela, on pourra démontrer que pour tout réel $x$ :

\[f(x) - x = \dfrac34(x - 2)^2.\]

\end{enumerate}

On considère la suite $\left(u_n\right)$ définie par un réel $u_0$ et pour tout entier naturel $n$ :
\[u_{n+1} = f\left(u_n\right).\]

On a donc, pour tout entier naturel $n$,

\[u_{n+1} = \dfrac34 u_n^2 - 2u_n + 3.\]

\begin{enumerate}[resume]
\item Étude du cas: $\dfrac43 \leqslant u_0 \leqslant  2$.
	\begin{enumerate}
		\item Démontrer par récurrence que, pour tout entier naturel $n$,
\[u_n \leqslant u_{n+1} \leqslant 2.\]

		\item En déduire que la suite $\left(u_n\right)$ est convergente.
		\item Prouver que la limite de la suite est égale à 2.
	\end{enumerate}
\item Étude du cas particulier : $u_0 = 3$.

On admet que dans ce cas la suite $\left(u_n\right)$ tend vers $+\infty$.

Recopier et compléter la fonction \og \texttt{seuil} \fg{} suivante écrite en Python, afin qu'elle renvoie la plus petite valeur de $n$ telle que $u_n$ soit supérieur ou égal à 100.

\begin{center}
\texttt{
\begin{tabular}{|l|}\hline
def seuil() :\\
\quad  u = 3\\
\quad  n = 0\\
\quad while \ldots\\
\quad\qquad  u = \ldots\\
\quad\qquad  n = \ldots\\
\quad return n\\ \hline
\end{tabular}}
\end{center}

\item Étude du cas : $u_0 > 2$.

À l'aide d'un raisonnement par l'absurde, montrer que $\left(u_n\right)$ n'est pas convergente.
\end{enumerate}

\bigskip
